\section{Evaluación humana multidimensional de Med-PaLM}

Una vez armada la cadena de modelos, la clase vuelve a la pregunta más importante sobre la evidencia: ¿Med-PaLM solo es más fuerte en preguntas de opción múltiple, o también ha reducido la distancia con los médicos en el rigor científico, el riesgo y la utilidad de las respuestas largas reales? Las nueve diapositivas de resultados que siguen no deben leerse solo por la altura de las barras o los colores, sino eje por eje.

\subsection{Consenso científico y evidencia de razonamiento}

La diapositiva de Scientific consensus compara Flan-PaLM, Med-PaLM y las clinician answers. La tendencia general es que Med-PaLM mejora de forma notable y se acerca a los expertos, aunque todavía queda una brecha. La siguiente diapositiva separa comprehension, knowledge retrieval y reasoning; en la sesión de preguntas se señaló en particular que la evidencia correcta y la evidencia incorrecta no son complementarias en el sentido de $1-x$: una misma respuesta puede contener al mismo tiempo hechos correctos e inferencias erróneas.

\lecturefigure{slide-35-clinical-consensus.jpg}{Concordancia con el consenso científico y clínico}{Diapositiva V035 recuperada de la grabación oficial; Stanford CS25 Lecture 17}

\lecturefigure{slide-36-comprehension-retrieval-reasoning.jpg}{Evidencia de comprensión, recuperación de conocimiento médico y razonamiento}{Diapositiva V036 recuperada de la grabación oficial; Stanford CS25 Lecture 17}

\teachervoice{Un estudiante preguntó en clase si los dos indicadores de comprehension eran complementarios, y el ponente respondió claramente «no». La comprensión correcta y la incorrecta pueden aparecer a la vez; por ejemplo, una respuesta puede identificar primero la enfermedad de forma correcta y luego añadir una inferencia irrelevante en la recomendación de tratamiento. Es una corrección muy importante para leer gráficas complejas de human-evaluation.}

\begin{knowledgebox}{Lectura de la figura: por qué los indicadores pueden ser verdaderos a la vez}
Una respuesta larga se compone de muchas proposition. Si una parte de ellas es correcta y otra parte es incorrecta, tanto «hay evidencia de comprensión correcta» como «hay evidencia de comprensión incorrecta» pueden ser verdaderas. Reducir todo el texto a correcto o incorrecto con un solo accuracy pierde este estado mixto; el valor de la evaluación eje por eje está precisamente en exponer las fallas locales.
\end{knowledgebox}

\subsection{Las respuestas más largas también pueden producir más errores}

La diapositiva de Incorrect or missing content revela un efecto secundario de la alineación: Med-PaLM se orienta a generar respuestas más completas y detalladas, pero el contenido adicional también amplía la oportunidad de que aparezcan afirmaciones irrelevantes o erróneas. El ponente explica que, en algunos ejes, el hecho de que Flan-PaLM parezca un poco mejor no significa que sea más seguro en conjunto; puede deberse simplemente a que dice menos. La evaluación debe considerar al mismo tiempo la completeness y la precision.

\lecturefigure{slide-37-incorrect-missing-context.jpg}{El doble riesgo del contenido erróneo y del contenido faltante}{Diapositiva V037 recuperada de la grabación oficial; Stanford CS25 Lecture 17}

\begin{importantbox}{La paradoja de la longitud en la generación clínica}
Una respuesta breve puede omitir advertencias críticas; una respuesta larga puede introducir alucinaciones adicionales. El objetivo no es maximizar el número de palabras, sino cubrir la información necesaria para la decisión mientras se controlan las proposition sin sustento. En términos de ingeniería, la «longitud de la respuesta» debe tratarse como una variable de riesgo, no como una señal natural de calidad.
\end{importantbox}

\subsection{Daño y sesgo: las consecuencias de los errores no son equivalentes}

La diapositiva de Harm considera a la vez la probabilidad y la gravedad. En la sesión de preguntas se dio una explicación concreta: el daño grave no proviene solo de recomendar un tratamiento peligroso, sino también de un misdiagnosis, de pasar por alto una condición que pone en riesgo la vida o de no comunicar la gravedad de la enfermedad. La diapositiva de Bias examina si la respuesta introduce sesgos demográficos inapropiados. Ambos aspectos requieren el juicio de expertos del dominio que tomen en cuenta el contexto; no pueden sustituirse con una lista negra de palabras clave.

\lecturefigure{slide-38-extent-likelihood-harm.jpg}{Grado y probabilidad del daño potencial}{Diapositiva V038 recuperada de la grabación oficial; Stanford CS25 Lecture 17}

\lecturefigure{slide-39-potential-bias.jpg}{Evaluación del sesgo potencial en las respuestas}{Diapositiva V039 recuperada de la grabación oficial; Stanford CS25 Lecture 17}

\begin{warningbox}{El daño no consiste en «si la respuesta contiene palabras peligrosas»}
Aunque el modelo no dé directamente una recomendación peligrosa, puede provocar un retraso en el tratamiento por un diagnóstico omitido, por subestimar la gravedad o por tranquilizar indebidamente. A la inversa, mencionar el cáncer o la muerte no constituye daño de forma automática; lo determinante es si el contenido corresponde a la pregunta y a la evidencia, y cómo podría cambiar la conducta del usuario.
\end{warningbox}

\teachervoice{El ponente concretó los ejemplos de harm grave como no identificar una patología mortal, por ejemplo un cancer, o no comunicar adecuadamente la gravedad. También señaló que el tema tiene muchos matices y que el artículo se apoya en el marco de riesgo de la AHRQ; la gráfica de la clase es solo una visión general de alto nivel y estas notas no deben reescribirla como una norma completa de seguridad clínica.}

\subsection{Lo que ve el usuario común es la helpfulness y la comprensibilidad}

La diapositiva de resultados de Layperson se centra en la intención y en el grado de ayuda. La tendencia aproximada que se presentó en clase es la siguiente: alrededor del 60\,\% de las respuestas de Flan-PaLM se consideran helpful, Med-PaLM sube a cerca del 80\,\% y las clinician answers se acercan al 90\,\%. Lo importante no es tomar estas proporciones como umbrales de despliegue, sino ver que la alineación médica puede reducir la brecha de comunicación, aunque sin alcanzar el nivel de los médicos.

\lecturefigure{slide-40-helpfulness-utility.jpg}{Evaluación de usuarios comunes: cobertura de la intención, grado de ayuda y utilidad}{Diapositiva V040 recuperada de la grabación oficial; Stanford CS25 Lecture 17}

\lecturefigure{slide-41-clinical-qualitative-example.jpg}{Caso cualitativo: diferencias entre las respuestas del médico, de Flan-PaLM y de Med-PaLM}{Diapositiva V041 recuperada de la grabación oficial; Stanford CS25 Lecture 17}

La diapositiva del ejemplo cualitativo debe leerse a partir de la forma en que se organiza la información: las respuestas del médico suelen ser precisas y breves, pero la terminología especializada puede dificultar que el paciente las entienda; el modelo puede desarrollar la explicación y añadir contexto, aunque al extenderse también puede introducir más errores. Lo que ambos se complementan es la conversión del lenguaje y la organización de la información, no que el modelo sustituya el juicio médico sin supervisión.

\begin{knowledgebox}{Lectura de la figura: cuatro aspectos que revisar primero al comparar respuestas}
Primero, verifique si la respuesta atiende directamente la intención del usuario; después, si están completos los hechos médicos esenciales; luego, si los detalles añadidos tienen evidencia, y por último, si se indica con claridad el límite a partir del cual hay que acudir a un profesional. La fluidez solo puede considerarse después de estos cuatro aspectos; de lo contrario, la respuesta mejor redactada podría juzgarse erróneamente como la más segura.
\end{knowledgebox}

El caso cualitativo también muestra los puntos ciegos de las metric automáticas. Dos respuestas pueden compartir gran parte del vocabulario y, sin embargo, tener consecuencias completamente distintas por una negación crucial, una dosis o una recomendación de acudir al médico; a la inversa, un médico puede dar la conclusión correcta con una sola abreviatura especializada y la similitud lingüística resultar muy baja. Una evaluación adecuada debe descomponer primero la respuesta en proposition clínicas, comprobar si cada una está respaldada por la pregunta y por el consenso, y después evaluar la organización, el tono y la posibilidad de actuar a partir de ella. Si el sistema se usa para reescribir las notas del médico para el paciente, también debe exigirse que conserve un enlace al texto original y un punto de confirmación humana, de modo que la «versión fácil de entender» no pueda alterar en silencio los hechos médicos. Así, la ventaja del modelo de lenguaje se aprovecha para explicar y convertir, mientras la fuente autorizada sigue siendo rastreable.

\subsection{¿Siguen haciendo falta los modelos de lenguaje clínicos?}

La penúltima diapositiva plantea una pregunta abierta: si los grandes modelos de propósito general ya son tan capaces, ¿siguen haciendo falta modelos de lenguaje especializados en lo clínico (clinical LM)? La clase no dio una respuesta negativa simple. Los modelos pequeños de dominio pueden aportar valor en costo, privacidad, despliegue local o distribuciones de datos específicas; los grandes modelos de propósito general, en cambio, aprovechan un conocimiento más amplio y el scaling. La verdadera pregunta es con qué datos, tareas y restricciones el beneficio marginal del preentrenamiento de dominio supera al de un modelo general más grande.

\lecturefigure{slide-42-clinical-lm-question.jpg}{Pregunta abierta: ¿siguen haciendo falta modelos de lenguaje clínicos especializados?}{Diapositiva V042 recuperada de la grabación oficial; Stanford CS25 Lecture 17}

\lecturefigure{slide-43-clinical-key-takeaways.jpg}{Conclusiones clave y limitaciones de la sección clínica}{Diapositiva V043 recuperada de la grabación oficial; Stanford CS25 Lecture 17}

\teachervoice{El papel que el ponente asigna a las aplicaciones clínicas a corto plazo no es el de un autonomous physician, sino el de augment physicians: convertir el doctor speak difícil de entender en explicaciones más accesibles para el paciente y ayudar a la comunicación entre médico y paciente, entre médicos y entre médicos e investigadores. Al mismo tiempo, dijo que esto todavía «no es perfecto» y que la oportunidad actual está en la complementarity.}

\begin{importantbox}{Conclusiones de la evidencia en la sección clínica}
La evidencia de la clase respalda tres puntos: los grandes modelos de propósito general transfieren con fuerza a las preguntas médicas de opción múltiple; el prompt tuning médico puede mejorar varios indicadores humanos de las respuestas largas, y el valor a corto plazo se parece más al apoyo a la comunicación y al conocimiento. No respalda que «aprobar un examen equivalga a obtener la habilitación clínica», ni que «una respuesta larga, por ser más detallada, sea más segura».
\end{importantbox}

\subsection{Resumen de la sección}

Med-PaLM supera claramente a Flan-PaLM en varias evaluaciones humanas, como el consenso científico, el razonamiento y el grado de ayuda, y reduce la brecha con las respuestas de los médicos; al mismo tiempo, las respuestas detalladas pueden introducir más errores, y el daño y el sesgo siguen requiriendo el juicio de expertos. La conclusión más confiable de la clase es el aumento clínico supervisado, no la sustitución automática.
